\chapter{Gauge Field Fluid Dynamics}
\label{chap:gauge-fluid}

% ── Brane fluid dynamics ────────────────────────────────────────────────────

\section{Fluid Dynamics of the Brane}
\label{sec:brane-fluid}

\begin{itemize}
  \item Describe the fluid dynamics on the brane whose thickness is
        determined by the substrate parameter~$q$
        (Definition~\ref{def:substrate-parameter}).
  \item TODO: investigate whether the viscosity gradient can be encoded
        by the $q$-Laplacian $\Delta_q$.
  \item TODO: formulate a Navier--Stokes analogue in the $q$-deformed
        gauge-field setting.
\end{itemize}

% ── Madelung remark ─────────────────────────────────────────────────────────

\begin{remark}[Madelung transformation: quantum fluid, geodesics, and divisor topology]
\label{rem:madelung-fluid}
The $q$-deformed Schr\"odinger equation
(Remark~\ref{rem:schrodinger-jet-coaction}) admits a fluid-dynamical
reformulation via the \emph{Madelung transformation}.  Writing the
linearised section $\psi \in \Gamma(E_{\mathrm{lin}})$ in polar form
\[
  \psi \;=\; \sqrt{\rho}\;\exp\!\bigl(J\,S/\hbar_q\bigr),
  \qquad
  \rho \;=\; \lvert\psi\rvert^2, \quad S \in C^\infty(M),
\]
and substituting into the $q$-deformed Schr\"odinger equation, one
obtains two real equations that are the direct analogues of classical
compressible fluid dynamics:
\begin{enumerate}
  \item \textbf{Continuity equation} (conservation of fluid density):
    \[
      \partial_t\,\rho \;+\; \nabla_{\!A} \cdot (\rho\,v) \;=\; 0,
      \qquad
      v \;:=\; \nabla_{\!A} S,
    \]
    expressing that $\rho = \lvert\psi\rvert^2$ is a conserved fluid
    density advected by the velocity field $v$.
  \item \textbf{Quantum Euler equation} (force balance):
    \[
      \partial_t v + (v \cdot \nabla_{\!A})\,v
      \;=\; -\nabla_{\!A}\!\bigl(V_q + Q_q\bigr),
      \qquad
      Q_q \;:=\; -\frac{\hbar_q^2}{2}\,
                  \frac{\Delta_q\sqrt{\rho}}{\sqrt{\rho}},
    \]
    where $Q_q$ is the \emph{$q$-deformed quantum pressure} (Bohm
    potential), the only term distinguishing the quantum fluid from a
    classical inviscid one.
\end{enumerate}

\medskip
\noindent\textbf{Fluid density and the substrate.}\;
The density $\rho = \lvert\psi\rvert^2$ is the actual density of the
gauge fluid at each point of the brane.  Via the substrate parameter
$q(p) = 1 + \lvert\langle L_\theta, F_A^+\rangle\rvert(p)$
(Remark~\ref{rem:c-substrate-via-levi}), the local value of $q$
modulates the \emph{effective thickness} of the brane fluid: where
$q(p) \approx 1$ (classical regime), the quantum pressure $Q_q$
vanishes and the fluid flows as an ordinary inviscid medium; where
$q(p) \gg 1$ (deep-quantum regime), $Q_q$ dominates and the fluid
resists compression, realising the inertial resistance identified with
mass (Remark~\ref{rem:mass-geometric-drag}).

\medskip
\noindent\textbf{The zero-information state: $\lvert\psi\rvert^2 = 0$.}\;
When $\rho(p) = \lvert\psi(p)\rvert^2 = 0$, the fluid density vanishes:
there is no medium at $p$ to carry a vortex, a signal, or a fermionic
excitation.  In the algebraic filter
(Remark~\ref{rem:algebraic-filter}), $p$ lies outside the support of
every harmonic mode, so no particle
(Definition~\ref{def:q-harmonic-form}) can be localised there.  This is
the geometric realisation of \emph{empty space}---not a fixed background
but a dynamically determined locus where the fibre of $E$ carries no
$\mathrm{SU}(2)$ content.  As $q(p) \to 1$ at such a point,
$\hbar_q(p) \to 0$ and quantum effects disappear entirely, consistent
with the vacuum being the state of maximum classical certainty and minimum
information.

\medskip
\noindent\textbf{Acoustic metric and analogue gravity.}\;
The two fluid equations above encode an \emph{acoustic metric}: small
perturbations $\delta\psi$ of the fluid propagate as sound waves and
``see'' only the background flow $v = \nabla_{\!A} S$, not the
underlying CR geometry directly.  Concretely, linearising the quantum
Euler equation around a stationary solution $(\rho_0, v_0)$ yields a
wave equation
\[
  \Box_{g_{\mathrm{ac}}}\,\delta\psi \;=\; 0,
\]
where $g_{\mathrm{ac}}$ is the \emph{acoustic metric}
\[
  g_{\mathrm{ac}}
  \;=\; \frac{\rho_0}{c_s}
  \begin{pmatrix}
    -(c_s^2 - \lvert v_0\rvert^2) & -v_0^{\,\top} \\
    -v_0 & \mathbf{1}
  \end{pmatrix},
  \qquad
  c_s^2 \;=\; \frac{\hbar_q^2}{4\rho_0}\,
              \Delta_q\rho_0,
\]
with $c_s$ the local speed of sound in the gauge fluid.  Where
$\lvert v_0\rvert = c_s$, the sound cone closes and the perturbation is
trapped: this is an \emph{acoustic horizon} (analogue black hole).

In the QOU framework (Remark~\ref{rem:gr-substrate-geometry}), General
Relativity is realised as the ``substrate geometry'': the Webster Ricci
curvature of $M$ equals the $q$-weighted stress-energy of the
$\mathrm{SU}(2)$ instantons.  The acoustic metric $g_{\mathrm{ac}}$
identifies this curvature with the non-uniform flow of the gauge fluid
--- regions of high torsion ($q \gg 1$) correspond to fast-moving
fluid, and acoustic horizons in the brane fluid are the geometric
realisation of gravitational horizons.  Photons (bosonic modes) travel
on null geodesics of $g_{\mathrm{ac}}$; fermions (matter modes) are
dragged into the flow and spiral around vortex filaments
(Remark~\ref{rem:algebraic-filter}).

\medskip
\noindent\textbf{Geodesics as stable vortex filaments.}\;
In the fluid-dynamical picture the velocity field $v = \nabla_{\!A} S$
is irrotational away from the zeros of $\psi$.  A trajectory along which
the quantum pressure $Q_q$ exerts no net transverse force---i.e.\
$(\nabla_{\!A} Q_q)_\perp = 0$ along the path---is a geodesic of the
acoustic metric $g_{\mathrm{ac}}$.  The algebraic filter
(Remark~\ref{rem:algebraic-filter}) then distinguishes two classes:
\begin{itemize}
  \item \emph{Bosonic geodesics} $([F_A, \alpha] = 0)$: the
        longitudinal modes are irrotational and travel straight through
        the bulk, following null geodesics of $g_{\mathrm{ac}}$.  These
        are the ``sound rays'' of the acoustic analogy.
  \item \emph{Fermionic geodesics} $([F_A, \alpha] \neq 0)$: the
        ladder modes are entangled with the curvature; the quantum
        pressure deflects them into closed orbits spiralling around the
        zeros of $\psi$.  Each such closed orbit is a \emph{stable
        vortex filament} $\gamma$, and its stability is controlled by the
        non-vanishing of the mass integral
        $m \propto \oint_\gamma \lVert\nabla_{\!A} q\rVert \cdot
        \lVert L \wedge F^+\rVert$
        (Definition~\ref{def:mass-identity}).
\end{itemize}

\medskip
\noindent\textbf{Linking and splitting in the divisor.}\;
The zero-locus $\{\psi = 0\}$ carries the structure of a (Weil)
\emph{divisor}: each irreducible component records a place where the
fluid empties out.  A topological blowup
(Conjecture~\ref{conj:topological-decomposition-mass}) replaces each
zero of $\psi$ by an exceptional divisor
$E_i \cong \mathbb{CP}^1$, and the quantum dilogarithm $\Phi_q(z_i)$
integrates the $q$-deformed Reeb flow along $E_i$, contributing a
localised mass increment to the vortex filament $\gamma$.  Two
operations on the divisor have direct fluid-dynamical meaning:
\begin{itemize}
  \item \emph{Splitting}: a single zero of $\psi$ resolves into two or
        more exceptional divisors $E_i$.  In the fluid picture, the
        vortex filament $\gamma$ \emph{bifurcates}: the density $\rho$
        first pinches off at the blowup centre, then the two sheets of
        $E_i$ carry independent vortex flows that diverge.  The mass
        decomposes as a sum $\sum_i \lambda_i\lvert\Phi_q(z_i)\rvert$
        over the contributions from each sheet
        (Conjecture~\ref{conj:topological-decomposition-mass}).
  \item \emph{Linking}: two exceptional divisors $E_i$, $E_j$ are
        \emph{linked} when the corresponding filaments $\gamma_i$,
        $\gamma_j$ form a non-trivial link in $M$---for instance, the
        Hopf link $\mathcal{L}_{\mathrm{Hopf}}$ associated with the
        strong force
        (Chapter~\ref{chap:path-integrals-braiding}).  Linking is
        detected by the intersection pairing $E_i \cdot E_j \neq 0$ in
        $H_*(M, \mathbb{Z})$.  In the fluid picture, the velocity fields
        of $\gamma_i$ and $\gamma_j$ wind around each other, producing a
        mutual induction whose energy is computed from the $q$-deformed
        skein relation (Chapter~\ref{chap:experimental}) and is the
        geometric origin of the strong-force binding energy.
\end{itemize}
Crucially, both operations are invisible at the level of $\rho =
\lvert\psi\rvert^2$ alone: they are topological properties of the
zero-\emph{locus} of $\psi$, not of the value of $\lvert\psi\rvert^2$.
The divisor structure---rather than the density itself---is therefore
the correct algebraic object for tracking vortex topology in the brane
fluid.
\end{remark}

% ── Geodesics and vortices ──────────────────────────────────────────────────

\section{Geodesics as Stable Vortex Points}
\label{sec:geodesic-vortices}

\begin{itemize}
  \item Geodesics on $M$ as stable fixed points of vortex configurations
        in the brane fluid.
  \item TODO: perform stability analysis via the $q$-deformed energy
        functional.
\end{itemize}

% ── Braid operators in fluid dynamics ───────────────────────────────────────

\section{Braid Operators in the Fluid Picture}
\label{sec:braid-fluid}

% TODO: work out the connection between braid-group operators
% (Chapter~\ref{chap:path-integrals-braiding}) and vortex braiding in the
% fluid dynamics.

% ── Planck limit and Reynolds number ────────────────────────────────────────

\section{$q$-Parameter, Planck Limit, and Reynolds Number}
\label{sec:planck-reynolds}

% TODO: relate the $q$-deformation parameter and the Planck limit to the
% critical separation distance and the Reynolds number of the brane fluid.
